\section{Decision Transformer: cómo una trayectoria se convierte en una token sequence}

Con un offline dataset, la trayectoria completa ya tiene registrada la reward futura en el momento del entrenamiento; por eso se puede calcular hacia atrás el return-to-go de cada instante e intercalarlo con el state y la action. Lo esencial es entender lo siguiente: durante el entrenamiento, el return-to-go proviene de trayectorias ya completadas; durante el despliegue, el usuario o el sistema fija un objetivo, que después se va descontando conforme llega la reward real.

\subsection{El return-to-go es una condición de objetivo, no «ver el futuro»}

En tareas de horizonte finito, el return-to-go sin descuento se define como
\[
\widehat{R}_t=\sum_{t'=t}^{T}r_{t'}.
\]
Si se requiere descuento, puede usarse
\[
\widehat{R}^{(\gamma)}_t=\sum_{t'=t}^{T}\gamma^{t'-t}r_{t'}.
\]
El símbolo $t$ denota el timestep actual, $T$ el final del episode, $r_{t'}$ la reward real del paso futuro $t'$ y $\gamma$ el discount factor. En la fase de entrenamiento, esta cantidad puede calcularse a partir de los registros; en la fase de inferencia solo se puede especificar un desired return y actualizar el objetivo restante según la reward ya obtenida.

\begin{importantbox}{Diferencia entre el return-to-go y la reward inmediata}
La reward inmediata $r_t$ responde «qué acaba de ocurrir»; el return-to-go $\widehat{R}_t$ responde «cuánto retorno total obtendrá todavía esta trayectoria desde ahora hasta el final». Decision Transformer usa este último como condición, de modo que un mismo modelo puede generar distintas distribuciones de acciones según el retorno objetivo.
\end{importantbox}

\subsection{Tres tipos de token y el causal Transformer}

Una trayectoria de longitud $T$ se ordena como
\[
\tau=(\widehat{R}_1,s_1,a_1,\widehat{R}_2,s_2,a_2,\ldots,\widehat{R}_T,s_T,a_T).
\]
Cada timestep aporta tres tokens: return, state y action. Cada uno pasa primero por su propio embedding, se suma después a un timestep embedding común y el resultado entra en el causal Transformer. El modelo produce $\widehat{a}_t$ en la posición de la action; la causal mask garantiza que la predicción solo use información actual y anterior.

\lecturefigure{slide-08-decision-transformer-architecture.jpg}{Arquitectura de Decision Transformer: return, state y action entran intercalados y el causal Transformer decodifica la acción en la posición de la action.}{Video oficial de Stanford Online, 12:56; corresponde a Chen et al. 2021.}

\begin{knowledgebox}{Lectura de la figura: seguir la información en una action prediction}
Empiece por $\widehat{R}_t,s_t,a_t$ en la parte inferior. Los bloques azules, verdes y naranjas representan los embeddings de distintos token types; el causal Transformer, situado encima de las posiciones, agrega un historial de longitud $K$; en la parte superior, el linear decoder predice la acción solo sobre los hidden states necesarios. En la figura, $a_t$ ya aparece como token del historial en posiciones posteriores, pero al predecir el $a_t$ actual el modelo no puede verlo a sí mismo; esto lo garantizan el desfase de la entrada y la causal mask.
\end{knowledgebox}

\begin{warningbox}{No tome el target return como una garantía}
Introducir $\widehat{R}_1=100$ solo expresa «se desea generar acciones coherentes con trayectorias de alto retorno»; no significa que el entorno vaya a devolver 100 al final. El modelo puede carecer de datos correspondientes, enfrentar transitions aleatorias o acumular errores en lazo cerrado. El requested return y el realized return deben registrarse por separado.
\end{warningbox}

\subsection{Perspectiva probabilística del conditional sequence modeling}

Desde el punto de vista de la aproximación de funciones, el modelo aprende
\[
p_{\theta}(a_t\mid \widehat{R}_{\le t},s_{\le t},a_{<t}).
\]
El parámetro $\theta$ representa todos los parámetros entrenables del Transformer y del embedding/decoder. Frente a una Markov policy que solo observa el state actual, esta distribución condicional conserva el historial de forma explícita; frente a la behavior cloning convencional, además toma el desired return como condición, lo que permite distinguir comportamientos de distinta calidad en torno a un mismo estado.

\begin{knowledgebox}{Términos clave: cuatro conceptos centrales de esta sección}
\begin{tabularx}{\textwidth}{L{3.0cm} X X}
\toprule
Término & Mecanismo & Malentendido frecuente \\
\midrule
causal mask & Impide que una posición acceda a tokens futuros & No equivale a que el modelo solo vea el state actual \\
timestep embedding & Marca la posición temporal dentro de la trayectoria & Resuelve un problema distinto del token-type embedding \\
context length $K$ & Ventana de historial que se conserva en cada predicción & Más larga no es necesariamente mejor; aumenta el cómputo y la necesidad de datos \\
conditional generation & Usa el retorno objetivo para elegir la distribución de comportamiento & No es una prueba de alcanzabilidad ni un task identifier universal \\
\bottomrule
\end{tabularx}
\end{knowledgebox}

\subsection{Resumen de la sección}

Decision Transformer intercala return, state y action en una causal sequence y aprende a predecir la action condicionada al historial y al retorno objetivo. El return-to-go se calcula hacia atrás a partir de los registros durante el entrenamiento y, durante el despliegue, es una condición de objetivo que se puede actualizar; ofrece una interfaz de control, pero no garantiza que el objetivo sea alcanzable.
